\section*{Contoh}
\addcontentsline{toc}{section}{Contoh}

\begin{contoh}[Menghitung closure]
Diberikan item \([E \to \cdot T]\) dan aturan
\(T \to T * F \mid F\). Closure menambahkan
\([T \to \cdot T * F]\) dan \([T \to \cdot F]\). Karena \(F\)
adalah nonterminal dengan aturan \(F \to (E) \mid \code{id}\),
closure juga menambahkan \([F \to \cdot (E)]\) dan
\([F \to \cdot \code{id}]\).
\end{contoh}

\begin{contoh}[Menghitung goto]
Untuk state yang memuat \([E \to E \cdot + T]\), goto pada simbol
\(+\) menghasilkan item \([E \to E + \cdot T]\), lalu closure
menambahkan aturan untuk \(T\). State baru ini kemudian memiliki
transisi pada simbol-simbol yang dapat mengikuti \(T\).
\end{contoh}

\begin{contoh}[Mengisi tabel SLR]
State yang memuat \([E \to E + T \cdot]\) menghasilkan aksi reduce
\(E \to E + T\) pada semua token di FOLLOW(\(E\)), yaitu
\(\{\$, ), +\}\). State yang memuat \([T \to T \cdot * F]\)
menghasilkan aksi shift pada \(*\).
\end{contoh}

\begin{contoh}[Konflik dangling else]
Pada state setelah \code{if E then stmt}, token \code{else} dapat
memicu shift (memasangkan else dengan if dalam) atau reduce
(mengurangi if luar). Tanpa aturan prioritas, tabel memiliki dua
aksi pada sel yang sama. Bison biasanya memilih shift, yang sesuai
dengan aturan else terdekat.
\end{contoh}

\begin{contoh}[Penggabungan LALR]
Dua state LR(1) dengan inti yang sama, misalnya
\(\{[A \to \alpha \cdot, a]\}\) dan
\(\{[A \to \alpha \cdot, b]\}\), digabung menjadi satu state LALR
\(\{[A \to \alpha \cdot, \{a,b\}]\}\). Penggabungan mengurangi
jumlah state, tetapi dapat menimbulkan konflik reduce-reduce yang
tidak ada pada LR(1) kanonik.
\end{contoh}
